\documentclass[12pt,a4paper,oneside]{article}
\usepackage[brazil]{babel}
\usepackage[utf8]{inputenc}
\usepackage{olymp}

\setcounter{problem}{3}

\begin{document}

\begin{problem}{Temperatura em Viçosa}{temp}{2}
 
A UFV recebe diversos estrangeiros para cursar disciplinas de graduação e pós-graduação. Muitos estranham as constantes mudanças bruscas de clima e de temperatura na cidade. Saem de casa com temperatura de $25^{\circ}$C e de repente uma chuva muda a temperatura para $15^{\circ}$C. Pior ainda para os americanos, que nem entendem essas temperaturas em $^{\circ}$C!

Faça um programa para converter temperaturas de $^{\circ}$C (graus Celsius) para $^{\circ}$F (graus Fahrenheit) e vice-versa. Para facilitar seu trabalho, damos a seguir a fórmula de conversão de $^{\circ}$F para $^{\circ}$C:

\[ \displaystyle ^{\circ}\text{C} = \frac{5}{9}(^{\circ}\text{F}-32) \]


%\Notes
%
%Observações, se tiver.

\InputFile

A entrada é composta por um valor real $T$ e um caractere $E$, sendo $T$ a temperatura e $E$ a escala utilizada, que pode ser `\texttt{F}' para Fahrenheit ou `\texttt{C}' para Celsius. Restrição: $-100 \leq T \leq 300$.


%\Notes

\OutputFile

Seu programa deve gerar apenas uma linha na saída, contendo um valor real, com uma casa decimal, representando a temperatura convertida, seguida por um espaço e um caractere (`\texttt{C}' ou `\texttt{F}') representando essa escala. Se a escala da entrada for Fahrenheit a temperatura na saída deve ser a correspondente em Celsius, e vice-versa.

\Example

\begin{example}
\exmp{
100 F
}{
37.8 C
}%
\end{example}

\begin{example}
\exmp{
100 C
}{
212.0 F
}%
\end{example}

%comment

\end{problem}

\end{document}
